% !TEX program = xelatex
% 2027 新领军学科能力二试 · 综合 · 化学（预测训练卷）
\documentclass[cjk=noto, seal=first]{mojing-exam}

\examtitle{2027 新领军学科能力测试 · 预测训练卷}
\secret{训练用卷 ★ 非官方试题}

\begin{document}

\begin{paper}[time={建议 1.5 小时（训练设置）}, full=100]{二试 · 综合 · 化学}

\begin{instructions}
  \item 官方只公布二试综合六选二的科目名称，没有公布化学的考纲、题型与分值。求真书院负责人在采访中表示，这一部分意在考查学生在数学之外
        是否怀有真实而持久的兴趣；本卷据此按“高中化学核心内容 + 定量推理”编写，全部安排为训练设置。
  \item 本卷共 4 题，每题 25 分，满分 100 分。可能用到的常数：$R=8.314\ \mathrm{J\,mol^{-1}\,K^{-1}}$，$F=96\,485\ \mathrm{C/mol}$，
        $\lg2\approx0.301$，$\ln2\approx0.693$，$\ln10\approx2.303$；Cu 的摩尔质量 $63.5\ \mathrm{g/mol}$。
\end{instructions}

\begin{gradingnotes}
  \item 本评分标准为训练用建议标准。数值结果在合理有效数字范围内即可；分值为累计得分。
\end{gradingnotes}

\scoretable{1,2,3,4}

\begin{questions}

% =============================================================================
\question[25]
（化学平衡）在一定温度下，$\mathrm{N_2O_4(g)\rightleftharpoons2NO_2(g)}$ 的标准平衡常数 $K^\ominus=\dfrac{(p_{\rm NO_2}/p^\ominus)^2}{p_{\rm N_2O_4}/p^\ominus}=0.15$，
$p^\ominus=100$ kPa。起始时容器中只有 $\mathrm{N_2O_4}$。
\begin{enumerate}
  \item 设平衡时 $\mathrm{N_2O_4}$ 的解离度为 $\alpha$、总压为 $p$，证明 $K^\ominus=\dfrac{4\alpha^2}{1-\alpha^2}\cdot\dfrac{p}{p^\ominus}$；
  \item 分别求总压为 $100$ kPa 和 $10$ kPa 时的 $\alpha$，并用勒夏特列原理解释其变化；
  \item 该反应正反应吸热。把装有平衡混合气的密闭玻璃球从冷水移入热水中，气体颜色怎样变化？$K^\ominus$ 怎样变化？
\end{enumerate}
\answerspace{2cm}
\begin{solution}
(1)~取 $1$ mol $\mathrm{N_2O_4}$，平衡时 $\mathrm{N_2O_4}$ 为 $1-\alpha$，$\mathrm{NO_2}$ 为 $2\alpha$，总量 $1+\alpha$。分压分别为
$\frac{1-\alpha}{1+\alpha}p$、$\frac{2\alpha}{1+\alpha}p$，代入得 $K^\ominus=\frac{4\alpha^2}{(1-\alpha)(1+\alpha)}\frac p{p^\ominus}$。\score{8}

(2)~$p=p^\ominus$：$\frac{4\alpha^2}{1-\alpha^2}=0.15$，$\alpha^2=\frac{0.15}{4.15}$，$\alpha\approx0.19$。
$p=0.1p^\ominus$：$\frac{0.4\alpha^2}{1-\alpha^2}=0.15$，$\alpha^2=\frac{0.15}{0.55}$，$\alpha\approx0.52$。\score{16}
减压时平衡向气体分子数增多的方向（解离方向）移动，解离度增大。\score{19}

(3)~吸热反应升温时 $K^\ominus$ 增大，平衡向生成红棕色 $\mathrm{NO_2}$ 的方向移动，颜色变深。\score{25}
\end{solution}

% =============================================================================
\question[25]
（缓冲溶液）25 ℃ 时醋酸的电离常数 $K_a=1.8\times10^{-5}$（$\mathrm pK_a\approx4.74$）。
\begin{enumerate}
  \item 求 $0.10\ \mathrm{mol/L}$ 醋酸溶液的 pH；
  \item 由 $K_a$ 的定义推导：含弱酸 HA 与其共轭碱 $\mathrm{A^-}$ 的溶液（二者浓度都远大于 $[\mathrm{H^+}]$）满足 $\mathrm{pH}=\mathrm pK_a+\lg\dfrac{[\mathrm{A^-}]}{[\mathrm{HA}]}$；
  \item 1.0 L 溶液含 0.10 mol HAc 与 0.10 mol NaAc，求其 pH；向其中加入 0.010 mol HCl（体积变化忽略），pH 变为多少？
        若把 0.010 mol HCl 加入 1.0 L 纯水中，pH 又是多少？由此说明缓冲溶液的作用及其限度。
\end{enumerate}
\answerspace{2cm}
\begin{solution}
(1)~$[\mathrm{H^+}]\approx\sqrt{K_ac}=\sqrt{1.8\times10^{-6}}\approx1.34\times10^{-3}$，pH $\approx2.87$。\score{6}

(2)~$K_a=\frac{[\mathrm{H^+}][\mathrm{A^-}]}{[\mathrm{HA}]}$，取负对数：$-\lg[\mathrm{H^+}]=-\lg K_a+\lg\frac{[\mathrm{A^-}]}{[\mathrm{HA}]}$。\score{12}

(3)~$\mathrm{pH}=\mathrm pK_a=4.74$。加酸后 $\mathrm{Ac^-}$ 变为 0.09 mol，HAc 变为 0.11 mol，pH $=4.74+\lg\frac{0.09}{0.11}\approx4.66$，
只下降约 0.09。\score{19}
纯水中 $[\mathrm{H^+}]=0.010$ mol/L，pH 由 7 变为 2。\score{21}
缓冲溶液中的共轭碱“吸收”加入的 $\mathrm{H^+}$（共轭酸吸收 $\mathrm{OH^-}$），使比值 $[\mathrm{A^-}]/[\mathrm{HA}]$ 变化不大；
但加入的酸或碱的量接近缓冲对本身的量时，缓冲作用失效。\score{25}
\end{solution}

% =============================================================================
\question[25]
（电化学）铜锌原电池 $\mathrm{Zn\,|\,Zn^{2+}\,\|\,Cu^{2+}\,|\,Cu}$ 的标准电动势 $E^\ominus=1.10$ V。25 ℃ 时能斯特方程为
$E=E^\ominus-\dfrac{0.0592}{n}\lg Q$，其中 $Q=\dfrac{[\mathrm{Zn^{2+}}]}{[\mathrm{Cu^{2+}}]}$，$n$ 为转移的电子数。
\begin{enumerate}
  \item 写出电极反应与总反应，求 $[\mathrm{Zn^{2+}}]=1.0$ mol/L、$[\mathrm{Cu^{2+}}]=0.010$ mol/L 时的电动势；
  \item 由 $\Delta G^\ominus=-nFE^\ominus$ 求标准反应吉布斯自由能，并求平衡常数 $K$ 的数量级（平衡时 $E=0$）；
  \item 用 2.0 A 的电流电解 $\mathrm{CuSO_4}$ 溶液 1.0 h，阴极析出铜多少克？
\end{enumerate}
\answerspace{2cm}
\begin{solution}
(1)~负极 $\mathrm{Zn-2e^-=Zn^{2+}}$，正极 $\mathrm{Cu^{2+}+2e^-=Cu}$，总反应 $\mathrm{Zn+Cu^{2+}=Zn^{2+}+Cu}$。
$E=1.10-\frac{0.0592}2\lg100\approx1.04$ V。\score{9}

(2)~$\Delta G^\ominus=-2\times96485\times1.10\ \mathrm J\approx-212$ kJ/mol。平衡时 $E=0$，$\lg K=\frac{2\times1.10}{0.0592}\approx37.2$，
$K\approx10^{37}$，反应进行得非常完全。\score{17}

(3)~$Q=It=7200$ C，$n(\mathrm e^-)=\frac{7200}{96485}$ mol，$m(\mathrm{Cu})=\frac12\times\frac{7200}{96485}\times63.5\approx2.37$ g。\score{25}
\end{solution}

% =============================================================================
\question[25]
（化学动力学）某物质在溶液中的分解是一级反应：$\dfrac{\dd c}{\dd t}=-kc$。300 K 时其半衰期为 20 min。
\begin{enumerate}
  \item 求速率常数 $k$，以及浓度降为初始值的 $10\%$ 所需的时间；
  \item 实验发现温度由 300 K 升到 310 K 时反应速率加倍。由阿伦尼乌斯公式 $k=A\ee^{-E_a/(RT)}$ 估算活化能 $E_a$；
  \item 加入催化剂后活化能降低 20 kJ/mol（指前因子 $A$ 不变）。300 K 时速率约变为原来的多少倍？说明“温度每升高 10 ℃ 速率约加倍”
        这一经验规则对应怎样的活化能范围。
\end{enumerate}
\answerspace{2cm}
\begin{solution}
(1)~$c=c_0\ee^{-kt}$，$k=\frac{\ln2}{20\ \mathrm{min}}\approx0.0347\ \mathrm{min^{-1}}$；$t=\frac{\ln10}{k}=20\times\frac{\ln10}{\ln2}\approx66$ min。\score{8}

(2)~$\ln2=\frac{E_a}R\Bigl(\frac1{300}-\frac1{310}\Bigr)$，$E_a=\frac{8.314\times0.693}{1.075\times10^{-4}}\approx5.4\times10^4$ J/mol $\approx54$ kJ/mol。\score{16}

(3)~$\frac{k'}{k}=\exp\frac{20\,000}{8.314\times300}\approx\ee^{8.0}\approx3\times10^3$。\score{21}
由 (2)，在室温附近“升温 10 ℃ 速率加倍”对应 $E_a$ 约为 50 kJ/mol 的数量级；活化能更大的反应对温度更敏感，规则只是粗略经验。\score{25}
\end{solution}

\end{questions}
\end{paper}
\end{document}
